\chapter{ANEXOS}

% Las tablas de los anexos no van al índice de tablas (corrección del centro de escritura).
\captionsetup{list=false}

\noindent
\begin{tabularx}{\textwidth}{@{}X r@{}}
  Documentación de la API REST & \pageref{anx:api} \\
  Comando de demostración (\texttt{seed\_demo}) & \pageref{anx:demo} \\
  Comandos reproducibles de los experimentos & \pageref{anx:comandos} \\
  Grilla de variantes y auditoría de robustez & \pageref{anx:grilla-variantes} \\
  Trazabilidad de requisitos al diseño & \pageref{anx:trazabilidad} \\
  Capturas complementarias de la interfaz & \pageref{anx:capturas} \\
  Modelo de datos de las bases legadas & \pageref{anx:modelo-legado} \\
\end{tabularx}

% Encabezado de cada anexo: salto de página, "Anexo N. Título" en negrita de subcapítulo, y entrada
% de índice general. #1 = número, #2 = título en la hoja del anexo, #3 = título en el índice general
% (igual a #2 salvo Anexo 2).
\newcommand{\anexo}[3]{%
  \clearpage
  \phantomsection
  \section*{\fontsize{12}{14}\selectfont\bfseries Anexo #1. #2}%
  \addcontentsline{toc}{section}{Anexo #1. #3}%
}

\anexo{1}{Documentación de la API REST}{Documentación de la API REST}
\label{anx:api}

La API se implementa con Django REST Framework y se publica bajo el prefijo \texttt{/api/}. La
autenticación es por \textit{JSON Web Token} mediante \textit{djangorestframework-simplejwt}: el
cliente obtiene un par de \textit{tokens} en \texttt{POST /api/auth/token/} y envía el
\textit{token} de acceso en la cabecera \texttt{Authorization: Bearer <token>}. La vigencia por
defecto es de 60 minutos para el \textit{token} de acceso y de 7 días para el de refresco, ambas
configurables por variable de entorno.

La autorización se resuelve sobre el campo \texttt{role} del usuario mediante dos permisos propios
del proyecto, \texttt{IsAdminRole} e \texttt{IsSurveyorRole}, y no sobre los permisos estándar de
Django REST Framework: \texttt{IsAdminUser} valida el atributo \texttt{is\_staff}, que en este
sistema no representa el rol de negocio. En las tablas siguientes, la columna Permiso distingue
cuatro niveles: Público (sin autenticación), Autenticado (cualquier usuario con sesión válida),
Admin (\texttt{IsAdminRole}) y Censista (\texttt{IsSurveyorRole}).

Los listados se paginan con \texttt{PageNumberPagination} a 20 elementos por página. Las respuestas
geográficas siguen el formato \textit{GeoJSON} (\textit{FeatureCollection}), con coordenadas en el
orden \texttt{[lon, lat]} que exige el estándar.

El inventario se reparte en cuatro tablas según el recurso: autenticación y usuarios
(Tabla~\ref{tab:anexo_api_auth}), \textit{datasets} e importación legada
(Tabla~\ref{tab:anexo_api_datasets}), motor de optimización (Tabla~\ref{tab:anexo_api_optimization})
y rutas y captura en terreno (Tabla~\ref{tab:anexo_api_routes}).

\begin{table}[H]
  \centering
  \caption{\textit{Endpoints} de autenticación y administración de usuarios.}
  \label{tab:anexo_api_auth}
  \small
  \begin{tabularx}{\textwidth}{@{}p{5.5cm} p{1.9cm} X@{}}
    \toprule
    Método y ruta &
      Permiso &
      Devuelve \\
    \midrule
    \texttt{POST \slash api\slash auth\slash token\slash } &
      Público &
      Par de \textit{tokens} JWT (\texttt{access}, \texttt{refresh}) a partir de usuario y
      contraseña. \\
    \texttt{POST \slash api\slash auth\slash token\slash refresh\slash } &
      Público &
      Nuevo \textit{token} de acceso a partir de un \textit{token} de refresco válido. \\
    \texttt{GET \slash api\slash auth\slash me\slash } &
      Autenticado &
      Usuario en sesión: \texttt{id}, \texttt{username}, \texttt{email} y \texttt{role}. Es la
      consulta que usa el \textit{frontend} para decidir qué interfaz mostrar. \\
    \texttt{GET \slash api\slash auth\slash surveyors\slash } &
      Admin &
      Censistas registrados, para poblar el selector de asignación de rutas. \\
    \texttt{GET \slash api\slash auth\slash users\slash } &
      Admin &
      Listado paginado de usuarios. \\
    \texttt{POST \slash api\slash auth\slash users\slash } &
      Admin &
      Crea un usuario. La contraseña es de solo escritura y se valida con las reglas de
      contraseña de Django. \\
    \texttt{GET \slash api\slash auth\slash users\slash \{id\}\slash } &
      Admin &
      Detalle de un usuario. \\
    \texttt{PATCH \slash api\slash auth\slash users\slash \{id\}\slash } &
      Admin &
      Actualiza rol, estado o contraseña. Rechaza que un administrador se degrade o desactive a
      sí mismo, y que se elimine el último administrador activo. \\
    \texttt{DELETE \slash api\slash auth\slash users\slash \{id\}\slash } &
      Admin &
      Desactivación lógica (\texttt{is\_active = false}); el usuario no se borra, para no
      arrastrar sus observaciones históricas. \\
    \bottomrule
  \end{tabularx}
\end{table}

\begin{table}[H]
  \centering
  \caption{\textit{Endpoints} de \textit{datasets} e importación legada.}
  \label{tab:anexo_api_datasets}
  \small
  \begin{tabularx}{\textwidth}{@{}p{5.5cm} p{1.9cm} X@{}}
    \toprule
    Método y ruta &
      Permiso &
      Devuelve \\
    \midrule
    \texttt{GET \slash api\slash datasets\slash } &
      Autenticado &
      Listado paginado de \textit{datasets} con su número de árboles. \\
    \texttt{POST \slash api\slash datasets\slash } &
      Admin &
      Crea un \textit{dataset} importando un archivo \texttt{.csv} o \texttt{.json}
      (\textit{multipart}). Detecta las columnas de coordenadas por nombre
      (\texttt{lat}/\texttt{latitude}/\texttt{latitud} y
      \texttt{lon}/\texttt{lng}/\texttt{longitude}/\texttt{longitud}) e infiere el delimitador
      del CSV. Devuelve el \textit{dataset} creado y \texttt{skipped\_rows}, el número de filas
      descartadas por coordenada faltante. \\
    \texttt{GET \slash api\slash datasets\slash \{id\}\slash } &
      Autenticado &
      Detalle de un \textit{dataset}. \\
    \texttt{PATCH \slash api\slash datasets\slash \{id\}\slash } &
      Admin &
      Actualiza nombre o descripción. \\
    \texttt{DELETE \slash api\slash datasets\slash \{id\}\slash } &
      Admin &
      Elimina el \textit{dataset} y, en cascada, sus árboles y soluciones. \\
    \texttt{GET \slash api\slash datasets\slash \{id\}\slash trees\slash } &
      Autenticado &
      Árboles activos del \textit{dataset} como \textit{FeatureCollection} de puntos. \\
    \texttt{GET \slash api\slash datasets\slash legacy\slash areas\slash } &
      Admin &
      Áreas de censo de la fuente legada con su campaña, su número de árboles y su polígono.
      Responde 503 si la fuente no está configurada. \\
    \texttt{GET \slash api\slash datasets\slash legacy\slash trees\slash } &
      Admin &
      Arbolado de ambas fuentes legadas con su posición, su especie, el área que lo contiene y
      la marca \texttt{already\_imported}. \\
    \texttt{POST \slash api\slash datasets\slash from-legacy-selection\slash } &
      Admin &
      Crea un \textit{dataset} en una única transacción a partir de una lista de árboles
      identificados por el par (\texttt{source}, \texttt{external\_id}), junto con su histórico
      de observaciones. \\
    \texttt{POST \slash api\slash datasets\slash import-legacy\slash } &
      Admin &
      Importa un área completa (\texttt{area\_id}) o la fuente entera (\texttt{area\_id} con
      valor \texttt{all}) sin selección previa. Atajo programático sin interfaz asociada. \\
    \texttt{GET \slash api\slash datasets\slash trees\slash \{id\}\slash observations\slash } &
      Autenticado &
      Historial completo de observaciones de un árbol, incluidas las importadas de las fuentes
      legadas. \\
    \bottomrule
  \end{tabularx}
\end{table}

\begin{table}[H]
  \centering
  \caption{\textit{Endpoints} del motor de optimización.}
  \label{tab:anexo_api_optimization}
  \small
  \begin{tabularx}{\textwidth}{@{}p{5.5cm} p{1.9cm} X@{}}
    \toprule
    Método y ruta &
      Permiso &
      Devuelve \\
    \midrule
    \texttt{GET \slash api\slash optimization\slash estimate\slash } &
      Admin &
      Estimación de flota para un \textit{dataset} (parámetros de consulta \texttt{dataset},
      \texttt{min\_route\_time\_sec} y \texttt{service\_time\_sec}). Se evalúa sobre la matriz
      de costos en caché y no dispara consultas a OSRM: si el \textit{dataset} no tiene matriz
      calculada, devuelve \texttt{n\_estimated} nulo. \\
    \texttt{POST \slash api\slash optimization\slash jobs\slash } &
      Admin &
      Crea la \texttt{RoutingConfig} ($T_{\min}$, $T_{\max}$, tiempo de servicio) y encola el
      \texttt{OptimizationJob} en Celery. El campo \texttt{strategy} admite \texttt{global},
      \texttt{spatial\_term}, \texttt{cluster\_first} o \texttt{compare} (ejecuta las tres y
      produce una solución por estrategia). \\
    \texttt{GET \slash api\slash optimization\slash jobs\slash } &
      Admin &
      Historial de trabajos, filtrable por \texttt{dataset}. \\
    \texttt{GET \slash api\slash optimization\slash jobs\slash \{id\}\slash } &
      Admin &
      Estado del trabajo (\texttt{queued}, \texttt{running}, \texttt{completed},
      \texttt{failed}), sus métricas, los árboles descartados y el mapa
      \texttt{solution\_ids} de estrategia a solución. Es el \textit{endpoint} que el
      \textit{frontend} consulta por \textit{polling}. \\
    \texttt{GET \slash api\slash optimization\slash solutions\slash \{id\}\slash } &
      Autenticado &
      Métricas de una solución: número de rutas, tiempo total de desplazamiento, puntaje de
      balance, métricas espaciales y tiempos por fase del \textit{pipeline}. El censista solo
      accede a las soluciones que contienen rutas suyas. \\
    \texttt{POST \slash api\slash optimization\slash solutions\slash \{id\}\slash publish\slash } &
      Admin &
      Publica la solución como plan vigente del \textit{dataset} y despublica atómicamente la
      anterior. \\
    \texttt{POST \slash api\slash optimization\slash solutions\slash \{id\}\slash
      unpublish\slash } &
      Admin &
      Retira la solución del plan vigente. \\
    \bottomrule
  \end{tabularx}
\end{table}

\begin{table}[H]
  \centering
  \caption{\textit{Endpoints} de rutas y captura en terreno.}
  \label{tab:anexo_api_routes}
  \small
  \begin{tabularx}{\textwidth}{@{}p{5.5cm} p{1.9cm} X@{}}
    \toprule
    Método y ruta &
      Permiso &
      Devuelve \\
    \midrule
    \texttt{GET \slash api\slash routes\slash } &
      Autenticado &
      Rutas con sus tiempos y su avance (visitadas, pendientes, omitidas), filtrables por
      \texttt{solution\_id}. El censista solo ve las rutas asignadas a él. \\
    \texttt{GET \slash api\slash routes\slash \{id\}\slash } &
      Autenticado &
      Detalle de una ruta, incluidas sus paradas ordenadas por \texttt{sequence}. \\
    \texttt{GET \slash api\slash routes\slash geojson\slash } &
      Autenticado &
      Polilíneas de todas las rutas de una solución (\texttt{solution\_id} obligatorio) como
      \textit{FeatureCollection} de \textit{LineString}, trazadas sobre la red peatonal real
      que devuelve OSRM. \\
    \texttt{GET \slash api\slash routes\slash \{id\}\slash path\slash } &
      Autenticado &
      Trazado peatonal de una sola ruta. \\
    \texttt{PATCH \slash api\slash routes\slash \{id\}\slash assign\slash } &
      Admin &
      Asigna o desasigna un censista. Solo se permite sobre rutas de la solución publicada. \\
    \texttt{GET \slash api\slash routes\slash my-route\slash } &
      Censista &
      Rutas propias del plan vigente; excluye por construcción las rutas de soluciones no
      publicadas. \\
    \texttt{POST \slash api\slash routes\slash stops\slash \{id\}\slash visit\slash } &
      Censista &
      Marca la parada como visitada y crea su \texttt{TreeObservation} (estado, fotografía y
      notas; \textit{multipart} cuando incluye imagen). Rechaza la visita si alguna parada
      anterior sigue pendiente. \\
    \texttt{POST \slash api\slash routes\slash stops\slash \{id\}\slash skip\slash } &
      Censista &
      Omite la parada con un motivo obligatorio y crea la observación correspondiente,
      derivando su estado del motivo. \\
    \bottomrule
  \end{tabularx}
\end{table}

\anexo{2}{Comando de demostración (\texttt{seed\_demo})}{Comando de demostración}
\label{anx:demo}

El comando de gestión \texttt{seed\_demo} construye un \textit{dataset} sintético y
reproducible sobre un rectángulo de Santiago
(latitud $[-33{,}45,\,-33{,}41]$; longitud $[-70{,}66,\,-70{,}58]$), de modo que la
plataforma pueda demostrarse y perfilarse sin depender de la disponibilidad de las
bases legadas. La generación es determinista: fijada la semilla
(\texttt{--seed}, $42$ por defecto), el mismo comando produce siempre el mismo
conjunto de puntos.

La Tabla~\ref{tab:anexo_seed_demo} detalla sus parámetros. Se ejecuta dentro del
contenedor \textit{backend}, que es donde residen las dependencias geográficas (GDAL)
y la conexión a PostGIS:

\begin{verbatim}
docker compose exec backend python manage.py seed_demo --profile light
docker compose exec backend python manage.py seed_demo \
    --trees 80 --distribution clustered --clusters 5 --snap
\end{verbatim}

\begin{table}[H]
  \centering
  \caption{Parámetros del comando \texttt{seed\_demo}.}
  \label{tab:anexo_seed_demo}
  \small
  \begin{tabularx}{\textwidth}{@{}p{5.5cm} X@{}}
    \toprule
    Parámetro &
      Efecto \\
    \midrule
    \texttt{--profile \{light, medium, heavy\}} &
      Perfiles predefinidos de carga: $15$, $70$ y $200$ árboles. \texttt{light} solo siembra
      datos; \texttt{medium} y \texttt{heavy} además resuelven el problema de ruteo. \\
    \texttt{--trees N} &
      Número de árboles cuando no se usa un perfil (por defecto $40$; mínimo $2$). \\
    \texttt{--distribution \{uniform, clustered, multizone\}} &
      Distribución espacial de los puntos. \texttt{uniform} muestrea el rectángulo completo;
      \texttt{clustered} genera nubes gaussianas de $250$\,m de dispersión en torno a centros
      aleatorios; \texttt{multizone} reparte los árboles entre tres zonas fijas de radios y
      pesos distintos, imitando la densidad desigual de un censo real. \\
    \texttt{--clusters K} &
      Número de centros de la distribución \texttt{clustered} (por defecto $4$). \\
    \texttt{--snap} &
      Proyecta cada punto generado al segmento de calle más cercano consultando
      \texttt{/nearest/v1/foot} de OSRM. Sin esta opción los árboles pueden caer en el interior
      de manzanas, lo que degrada los tiempos de la matriz de costos. \\
    \texttt{--seed S} &
      Semilla del generador pseudoaleatorio (por defecto $42$). \\
    \texttt{--name} &
      Nombre del \textit{dataset} creado (por defecto \textit{«Demo Santiago»}). \\
    \texttt{--optimize / --no-optimize} &
      Fuerza o inhibe la ejecución del \textit{pipeline} de optimización sobre el
      \textit{dataset} recién creado, con independencia del perfil. \\
    \bottomrule
  \end{tabularx}
\end{table}

El comando crea un \texttt{Dataset} y sus \texttt{Tree}, un \texttt{PointField} por
árbol con SRID 4326. Cuando el perfil o el \textit{flag} lo indican, crea además la
\texttt{RoutingConfig} con los parámetros por defecto ($T_{\min} = 7\,200$\,s,
$T_{\max} = 10\,800$\,s, tiempo de servicio $120$\,s), el \texttt{OptimizationJob}
correspondiente y ejecuta el \textit{pipeline} de forma síncrona ---sin pasar por
Celery---, lo que persiste la \texttt{RoutingSolution} con sus \texttt{Route} y
\texttt{RouteStop}. La ejecución síncrona es deliberada: permite reproducir una
corrida completa desde la línea de comandos sin levantar el \textit{worker}.

El \textit{dataset} de la evaluación de esta memoria no proviene de este comando sino
de las fuentes legadas (Sección~\ref{sec:legacy}); \texttt{seed\_demo} cubre la
demostración funcional y las pruebas de carga sintéticas.

\anexo{3}{Comandos reproducibles de los experimentos}{Comandos reproducibles de los experimentos}
\label{anx:comandos}

Los resultados del Capítulo~\ref{cap:resultados} se obtuvieron con tres comandos de
gestión: \texttt{baseline\_sweep} (barridos de configuración y perfilado por fases),
\texttt{baseline\_postpass} (corridas de referencia de \textit{OR-Tools} con
post-proceso habilitado) y \texttt{greedy\_baseline} (el heurístico de contraste, con
y sin post-proceso de resecuenciación). Las tablas siguientes detallan los parámetros
de cada uno. Cada ejecución escribe además un informe en \texttt{docs\slash
experiments\slash }, con su comando, sus parámetros y sus métricas.

La versión de la biblioteca de optimización queda fijada en
\texttt{backend\slash requirements.txt}: todas las corridas del
Capítulo~\ref{cap:resultados} emplean \texttt{ortools==9.15.6755}. Reproducir las cifras
exige esa versión, porque la biblioteca no expone una semilla y las heurísticas de
construcción y de búsqueda local cambian entre versiones (RE-05).

La Tabla~\ref{tab:anexo_baseline_sweep} detalla los parámetros de \texttt{baseline\_sweep}.

\begin{table}[H]
  \centering
  \caption{Parámetros del comando \texttt{baseline\_sweep}.}
  \label{tab:anexo_baseline_sweep}
  \small
  \begin{tabularx}{\textwidth}{@{}p{5.5cm} X@{}}
    \toprule
    Parámetro &
      Efecto \\
    \midrule
    \texttt{--dataset <uuid>} &
      Barre un \textit{dataset} real ya existente. Si se omite, el comando genera
      \textit{datasets} sintéticos según \texttt{--sizes} y \texttt{--distribution}. \\
    \texttt{--strategies} &
      Lista separada por comas: \texttt{global}, \texttt{spatial\_term},
      \texttt{cluster\_first}. \\
    \texttt{--service-time} &
      Lista de tiempos de servicio por árbol, en minutos. \\
    \texttt{--t-max} &
      Lista de duraciones máximas de ruta, en horas. \\
    \texttt{--seeds N} / \texttt{--base-seed} &
      Número de repeticiones y semilla inicial. En \texttt{baseline\_sweep} la semilla no llega
      al \textit{solver}: sobre un \textit{dataset} real solo etiqueta la repetición, y la
      varianza observable proviene del corte por tiempo de reloj de la búsqueda local (sobre
      \textit{datasets} sintéticos, además genera instancias distintas por semilla). \\
    \texttt{--time-limit <s>} &
      Fija el límite de tiempo del \textit{solver}, en lugar de la heurística del
      \textit{pipeline} ($\min(30 + 1{,}5\,n,\;120)$ segundos). \\
    \texttt{--csv <ruta>} &
      Ruta del CSV de salida. \\
    \bottomrule
  \end{tabularx}
\end{table}

El barrido de perfilado por fases ---seis variantes operacionales de tiempo de
servicio y $T_{\max}$ sobre el \textit{dataset} real, con la estrategia \texttt{global}
y tres repeticiones--- se reproduce con:

\begin{verbatim}
docker compose -f docker-compose.prod.yml run --rm backend \
  python manage.py baseline_sweep \
    --dataset <uuid> --strategies global \
    --service-time 1,2,3 --t-max 2,3 \
    --seeds 3 --time-limit 180 \
    --csv /results/profiling-grid-full.csv
\end{verbatim}

La comparación directa de las tres estrategias de partición, bajo una misma
variante operacional y un mismo presupuesto de \textit{solver}, se reproduce con:

\begin{verbatim}
docker compose -f docker-compose.prod.yml run --rm backend \
  python manage.py baseline_sweep \
    --dataset <uuid> \
    --strategies global,spatial_term,cluster_first \
    --service-time 2 --t-max 3 \
    --seeds 3 --base-seed 42 --time-limit 120 \
    --csv /results/strategy-comparison.csv
\end{verbatim}

\begin{sloppypar}
Ambos barridos se ejecutan sobre las imágenes de producción
(\texttt{docker-compose.prod.yml}) para que los tiempos medidos correspondan a esas
imágenes y no a las de desarrollo. El protocolo completo de cada experimento
---entorno, criterios de decisión fijados antes de correr, resultados y su análisis---
se documenta, en el directorio \texttt{docs\slash experiments\slash }, en los informes
\texttt{profiling-20260711.md} (perfilado por fases) y
\texttt{route-quality-20260712.md} (comparación de estrategias). Las cifras que
sustentan el Capítulo~\ref{cap:resultados} provienen de los CSV depositados en ese
mismo directorio.
\end{sloppypar}

El comando \texttt{baseline\_postpass} ejecuta la corrida de referencia de
\textit{OR-Tools} con el post-proceso de resecuenciación habilitado por defecto,
consolidando múltiples semillas en un CSV. La
Tabla~\ref{tab:anexo_baseline_postpass} detalla sus parámetros. Se ejecuta dentro
del contenedor \textit{backend}:

\begin{verbatim}
docker compose -f docker-compose.prod.yml run --rm backend \
  python manage.py baseline_postpass \
    --dataset reference-n1607 \
    --strategy spatial_term \
    --seeds 42,43,44
\end{verbatim}

\begin{sloppypar}
Los datos crudos de la corrida de referencia de la Sección~\ref{sec:metricas-calidad} ---una fila
por semilla--- son los que ese comando escribe en
\texttt{20260724-225420-postpass-baseline.csv}. Las cifras que esa misma sección toma de corridas
distintas provienen de \texttt{20260713-real-case-metrics-spatial.csv} (tiempo de reloj extremo a
extremo con la matriz en caché), de \texttt{reference-n1607-20260730.csv} (auto-cruces sobre el
trazado de calles) y de \texttt{penalty-sensitivity-20260713-reference.csv}, columna
\texttt{tmin\_gap\_routes} (rutas por debajo de $T_{\min}$). Los resultados de la corrida de
referencia anterior ---estrategia \texttt{global} y límite de $180$\,s, los \textit{defaults}
previos al ajuste de la Sección~\ref{sec:comparacion-estrategias}--- se conservan en
\texttt{20260712-real-case-metrics.csv}.
\end{sloppypar}

\begin{table}[H]
  \centering
  \caption{Parámetros del comando \texttt{baseline\_postpass}.}
  \label{tab:anexo_baseline_postpass}
  \small
  \begin{tabularx}{\textwidth}{@{}p{5.5cm} X@{}}
    \toprule
    Parámetro &
      Efecto \\
    \midrule
    \texttt{--dataset <nombre>} &
      Nombre de la instancia congelada (por defecto \texttt{reference-n1607}). \\
    \texttt{--strategy <str>} &
      Estrategia del \textit{solver}: \texttt{global}, \texttt{spatial\_term} o
      \texttt{cluster\_first} (por defecto \texttt{spatial\_term}). \\
    \texttt{--seeds <lista>} &
      Lista de semillas separadas por comas (por defecto \texttt{42,43,44}). Cada semilla
      genera una permutación de nodos que alimenta la heurística de construcción del GLS. \\
    \texttt{--csv <ruta>} &
      Ruta del CSV de salida. Si se omite, el comando escribe en \texttt{docs/experiments/}
      con marca temporal. \\
    \bottomrule
  \end{tabularx}
\end{table}

El comando \texttt{greedy\_baseline} ejecuta el heurístico de vecino más cercano
sobre un \textit{dataset} real, con o sin post-proceso de resecuenciación 2-opt, y
emite las mismas métricas de calidad que \texttt{baseline\_sweep} para comparación
directa. Es determinista: no utiliza semilla aleatoria. La
Tabla~\ref{tab:anexo_greedy_baseline} detalla sus parámetros.

\begin{verbatim}
docker compose -f docker-compose.prod.yml run --rm backend \
  python manage.py greedy_baseline \
    --dataset <uuid> \
    --service-time 2 --t-max 3 \
    --postpass
\end{verbatim}

\begin{table}[H]
  \centering
  \caption{Parámetros del comando \texttt{greedy\_baseline}.}
  \label{tab:anexo_greedy_baseline}
  \small
  \begin{tabularx}{\textwidth}{@{}p{5.5cm} X@{}}
    \toprule
    Parámetro &
      Efecto \\
    \midrule
    \texttt{--dataset <uuid>} &
      UUID del \textit{dataset} (obligatorio). \\
    \texttt{--service-time <min>} &
      Tiempo de servicio por árbol en minutos (por defecto $2$\,min). \\
    \texttt{--t-max <h>} &
      Cota dura $T_{\max}$ en horas (por defecto $3$\,h). \\
    \texttt{--postpass} &
      Aplica el mismo post-proceso de resecuenciación 2-opt que recibe la solución del
      \textit{solver}, de modo que ambos enfoques reciben el mismo tratamiento. \\
    \texttt{--csv <ruta>} &
      Ruta del CSV de salida. Si se omite, el comando escribe en \texttt{docs/experiments/}
      con marca temporal. \\
    \bottomrule
  \end{tabularx}
\end{table}

\anexo{4}{Grilla de variantes y auditoría de robustez}{Grilla de variantes y auditoría de robustez}
\label{anx:grilla-variantes}

Este anexo complementa las Secciones~\ref{sec:variantes-duracion} y~\ref{sec:discusion-calidad} con
el detalle de la grilla de variantes operacionales, la descripción de la suite de instancias
congeladas sobre la que se ejecutó el barrido de configuración y la descripción completa de los
cuatro defectos de instrumentación detectados en la serie experimental.

\subsubsection*{Grilla de seis variantes}

La Tabla~\ref{tab:grilla-variantes} reporta las seis combinaciones de tiempo de servicio por árbol
$\{1, 2, 3\}$\,min y cota dura $T_{\max} \in \{2, 3\}$\,h sobre el \textit{dataset} real de
$n=1\,607$ árboles (429 de \texttt{legacy\_api} y 1\,178 de \texttt{legacy\_app}, con posición
mediana por código QR). Cada variante se ejecutó con estrategia \texttt{global}, semillas 42, 43 y
44, un límite de \textit{solver} de 180\,s y la versión del código previa a la corrección del
redondeo del tránsito; la tabla reporta medias. El total del \textit{pipeline} equivale al tiempo de
\texttt{solve} más $\approx$0{,}6\,s de sobrecosto fijo.

\begin{table}[H]
  \centering
  \caption{Grilla de seis variantes sobre el \textit{dataset} real.}
  \label{tab:grilla-variantes}
  \small
  \begin{tabularx}{\textwidth}{rrXXXXXXX}
    \toprule
    Serv. [min] &
      $T_{\max}$ [h] &
      $k$ &
      Balance &
      $\bar{T}$ [s] &
      $\sigma(T)$ [s] &
      Descart. &
      Desplaz. [s] &
      \texttt{solve} [s] ($\pm$\,sd) \\
    \midrule
    1 & 2 & 23{,}0 & 0{,}992 & 7\,229  & 14  & 0     & 69\,868 & 92{,}8 $\pm$ 63{,}8 \\
    1 & 3 & 15{,}0 & 0{,}877 & 10\,417 & 378 & 0     & 59\,843 & 89{,}0 $\pm$ 3{,}6 \\
    2 & 2 & 36{,}0 & 0{,}996 & 7\,219  & 6   & 0     & 67\,047 & 131{,}6 $\pm$ 34{,}4 \\
    2 & 3 & 24{,}0 & 0{,}865 & 10\,399 & 391 & 0     & 56\,750 & 126{,}9 $\pm$ 6{,}7 \\
    3 & 2 & 49{,}0 & 0{,}998 & 7\,213  & 3   & 1{,}0 & 64\,360 & 139{,}2 $\pm$ 14{,}1 \\
    3 & 3 & 33{,}0 & 0{,}890 & 10\,554 & 284 & 0     & 59\,023 & 49{,}3 $\pm$ 14{,}1 \\
    \bottomrule
  \end{tabularx}
\end{table}

Los resultados muestran el compromiso central entre balance y costo total. Con $T_{\max}=2$\,h el
número de rutas $k$ escala fuertemente con el tiempo de servicio
($23 \rightarrow 36 \rightarrow 49$), dado que el trabajo de servicio domina la capacidad de cada
ruta, y el balance se aproxima a 1 (al coincidir $T_{\min}$ y $T_{\max}$, las rutas resultan de
duración casi idéntica, con $\sigma(T)$ de 3 a 14\,s). Con $T_{\max}=3$\,h, en cambio, el tiempo
total de desplazamiento se reduce en torno al 10\,\% y $k$ cae alrededor de un 35\,\%, a costa de un
balance de $\approx 0{,}87$ con $\sigma(T)$ de 284 a 391\,s. La única variante que descarta nodos es
la más restrictiva (servicio de 3\,min con $T_{\max}=2$\,h), que deja un árbol fuera de la solución
de forma consistente en las tres semillas.

Respecto de las semillas, corresponde precisar su rol: a diferencia de
\texttt{config\_algorithm\_sweep} (Sección~\ref{sec:discusion-calidad}), donde la semilla sí se
propaga al \textit{solver} y altera la trayectoria del GLS, esta grilla se ejecutó con
\texttt{baseline\_sweep} (\hyperref[anx:comandos]{Anexo 3}), donde la semilla solo etiqueta la
repetición. La varianza observada entre repeticiones (columna sd de \texttt{solve} en la tabla)
proviene íntegramente del corte por tiempo de reloj (\textit{wall-clock}) de la búsqueda local: en
las variantes en que la búsqueda agota el presupuesto, las tres repeticiones convergen a soluciones
casi idénticas.

\phantomsection
\subsubsection*{Suite de instancias congeladas}
\label{anx:suite-instancias}

Las doce instancias sobre las que se ejecutó el barrido de configuración
(Sección~\ref{sec:discusion-calidad}) están versionadas en \texttt{docs/experiments/instances/} y se
describen en la Tabla~\ref{tab:suite-instancias}. La batería de crecimiento se construye sobre los
$12\,946$ árboles georreferenciados de \texttt{legacy\_app} ---los $12\,941$ del importador más los
cinco códigos QR de ceros, que la batería no filtra por construirse antes de esa exclusión---,
ordenados por distancia de \textit{haversine} a un centroide fijo: cada instancia es prefijo exacto
de la siguiente, de modo que el crecimiento es puramente espacial y no una muestra distinta. Las dos
variantes dispersas toman uno de cada dos y uno de cada cuatro sobre los primeros $1\,000$, con lo
que conservan la extensión y reducen la densidad a la mitad y a un cuarto: aíslan el efecto de la
densidad del efecto del tamaño.

\begin{table}[H]
  \centering
  \caption{Suite de instancias congeladas empleada en el barrido de configuración.}
  \label{tab:suite-instancias}
  \small
  \begin{tabularx}{\textwidth}{@{}l r X@{}}
    \toprule
    Familia & $n$ & Papel en la serie \\
    \midrule
    \texttt{area-29-n43}, \texttt{area-27-n72}, \texttt{area-26-n157} &
      43--157 &
      Áreas reales del censo histórico: la unidad operativa. \\
    \texttt{battery-n50} a \texttt{battery-n1000} (seis instancias) &
      50--1\,000 &
      Batería de crecimiento anidada sobre \texttt{legacy\_app}: aísla el tamaño a
      densidad real. \\
    \texttt{battery-sparse-n250}, \texttt{battery-sparse-n500} &
      250, 500 &
      Misma extensión que \texttt{battery-n1000} con densidad a un cuarto y a la mitad:
      aísla la densidad del tamaño. \\
    \texttt{reference-n1607} &
      1\,607 &
      Régimen agregado: la unión de ambas fuentes legadas. \\
    \bottomrule
  \end{tabularx}
\end{table}

\phantomsection
\subsubsection*{Reservas del experimento de sensibilidad a penalizaciones}
\label{anx:reservas-sensibilidad}

La Sección~\ref{sec:discusion-calidad} resume en una frase las tres reservas que acotan la lectura
de la reducción de auto-cruces al redirigir la penalización blanda superior hacia $T_{\max}$
(\texttt{docs/experiments/penalty-sensitivity-20260713-reference.csv}). Este apartado las detalla.

La primera es de replicación. Las tres filas de cada configuración registran $76$, $77$ y $76$
auto-cruces en un brazo y $6$, $7$ y $6$ en el otro: coinciden a la unidad porque no son réplicas
independientes. Ese barrido es anterior a la corrección de instrumento del Defecto~1 más abajo, y el
comando de gestión que lo ejecutó escribía la semilla en el CSV sin propagarla al \textit{solver}.
Las cifras son, por tanto, una medición por celda y no la media de tres repeticiones, y la reducción
no lleva barra de error.

La segunda es de instrumento. La métrica de auto-cruces se calcula sobre cuerdas rectas entre
paradas consecutivas, y la auditoría de instrumento de julio de 2026
(\texttt{docs/experiments/evidence-audit-20260724.md}) midió su correlación con el trazado real de
calles: un $\rho$ de Spearman global de $0{,}527$, propio de un indicador indirecto parcial, e
\textit{invertido} en el \textit{dataset} de evaluación agregado ($\rho = -0{,}575$), que es
precisamente la instancia sobre la que se midió esta reducción. La cifra de $91{,}7$\,\% describe lo
que ocurre con las cuerdas, y en esta instancia la cuerda es un mal predictor del signo del cambio
en la calle.

La tercera es de costo omitido. El mismo CSV muestra que el brazo que redirige la penalización
superior a $T_{\max}$ degrada el balance de carga de $0{,}838$ a $0{,}715$ y eleva de tres a cinco
el número de rutas por debajo de $T_{\min}$. El balance es uno de los tres propósitos declarados del
modelo (Sección~\ref{sec:formulacion}), de modo que la ganancia geométrica no es gratuita: se paga
en la dimensión que el criterio de aceptación CA-03 protege.

Establecer la afirmación en su forma fuerte exigiría re-correr el experimento con semillas
efectivamente propagadas al \textit{solver} y reportar los auto-cruces sobre el trazado de calles,
instrumentación que el repositorio ya soporta y que este trabajo no ejecutó.

\phantomsection
\subsubsection*{Barrido de configuración $\times$ algoritmo $\times$ tamaño --- detalle}
\label{anx:barrido-configuraciones}

La Sección~\ref{sec:discusion-calidad} resume el barrido ejecutado sobre las doce instancias
congeladas de este anexo (\texttt{route-config-algorithm-sweep-20260718.csv}: ocho configuraciones
sobre cada instancia, 96 celdas, registradas en 288 filas). Cada celda se registró con tres
repeticiones nominales, pero el comando de gestión que ejecutó el barrido escribía la semilla en el
CSV sin propagarla al \textit{solver}, de modo que las tres filas de una misma celda no son réplicas
independientes: difieren solo por el corte por reloj de pared. En 78 de las 96 celdas las tres filas
coinciden exactamente; en las 18 restantes el desplazamiento total varía hasta un $3{,}4$\,\%. Las
cifras que siguen provienen, por tanto, de una medición por celda sin réplica de semilla y no llevan
barras de error.

El brazo \texttt{upper-tmax-tmin9000} ---que ancla la penalización blanda superior en $T_{\max}$ en
lugar del punto medio del intervalo $[T_{\min},\,T_{\max}]$ y eleva el piso blando a $9\,000$\,s, a
diferencia del brazo de sensibilidad anterior que solo redirigía la cota superior sin elevar el
piso--- es el único que satisface el criterio de decisión para $n=1\,607$: los auto-cruces caen un
$92{,}5$\,\% ($89{,}0 \to 6{,}7$) con una variación de viaje de $-0{,}3$\,\%, balance de $0{,}84$,
$k=25$ y cero descartes. La razón es estructural: a esa densidad las rutas saturan el presupuesto
temporal de forma natural ($\overline{\text{sat}} = 0{,}94$), de modo que anclar la cota superior en
$T_{\max}$ no impone \textit{detours} adicionales. En el eje de algoritmo, \texttt{global}
prácticamente no modifica los auto-cruces ($-2{,}6$\,\%) y \texttt{greedy} los empeora
($+14{,}6$\,\%), lo que ratifica \texttt{spatial\_term} como la estrategia correcta.

Sin embargo, el mismo brazo produce el peor resultado en las áreas reales individuales, donde la
instancia no satura $T_{\max}$ ($\overline{\text{sat}} \approx 0{,}67$--$0{,}73$): anclar la cota
superior fuerza relleno de trayecto porque la carga de servicio disponible no llena el presupuesto,
con incrementos de viaje de $+64{,}0$\,\%, $+62{,}6$\,\% y $+89{,}0$\,\% en las instancias de $n =
157$, $72$ y $43$, respectivamente. Bajo la configuración de producción, esas mismas áreas generan
cero o casi cero auto-cruces; el auto-cruce es un fenómeno exclusivo del régimen denso agregado. El
brazo \texttt{span-c100} ---penalización del \textit{span} sobre la dimensión temporal mediante
\texttt{SetSpanCostCoefficientForAllVehicles}, que eleva el coeficiente de $0$ en producción a $100$
y deja intacta la penalización de $3$ sobre la dimensión espacial--- confirma la predicción de
nulidad: solo reduce los auto-cruces en $26{,}6$\,\% para $n=1\,607$, por debajo del umbral del
$30$\,\% exigido por el criterio a priori, porque penalizar la suma de duraciones de ruta no
distingue el relleno del viaje productivo. Los brazos que reducen o suprimen el piso blando
(\texttt{tmin-scaled}, \texttt{service-floor} y \texttt{tmin-scaled+exempt-last}) quedan
descalificados como configuración universal al romper la compuerta de balance $\geq 0{,}80$ en
siete, siete y cuatro instancias, respectivamente.

La Sección~\ref{sec:posibles-mejoras} menciona una compuerta de régimen por saturación, descartada,
que activaría \texttt{upper-tmax-tmin9000} por encima de $\overline{\text{sat}} \gtrsim 0{,}9$. La
re-medición sobre \texttt{crossings\_road} invierte el signo del efecto precisamente en el régimen
que la compuerta habilitaría: la mediana de auto-cruces sobre el trazado de calles sube de $43$ a
$69$ en \texttt{reference-n1607} ($\overline{\text{sat}} = 0{,}94$), de $46$ a $75$ en
\texttt{battery-n1000} ($0{,}91$) y de $26$ a $59$ en \texttt{battery-n800} ($0{,}89$).

\phantomsection
\subsubsection*{Auditoría de robustez metodológica --- detalle}
\label{anx:auditoria-robustez}

La Sección~\ref{sec:discusion-calidad} resume los cuatro defectos de
instrumentación detectados en la serie experimental. Este apartado los describe
con el detalle necesario para reproducir las correcciones.

Defecto 1 --- Semillas no propagadas.
Las corridas iniciales de la serie no propagaban la semilla al \textit{solver}: las
nominales tres réplicas por celda eran la misma corrida, de modo que la dispersión
reportada ($\sigma \approx 0{,}03\,\%$) era un artefacto de replicación idéntica y no
describía la convergencia del algoritmo. Una corrida de recalibración posterior propagó
semillas reales al \textit{solver}, con lo que la dispersión subió a $\sigma \approx
2{,}4\,\%$ y catorce de veintiocho veredictos de ese ciclo debieron re-juzgarse.

Defecto 2 --- Cota de relleno inalcanzable.
Los ciclos iniciales usaban como piso de la métrica de relleno el producto
$n \times \bar{n}$ (número de nodos por grado medio de la red), que viola la restricción
de grado de los caminos y resulta por tanto inalcanzable. La métrica corregida usa el
bosque generador mínimo de $k$ componentes ($\text{MSF}_k$) \citep{kruskal1956shortest}
sobre la matriz simetrizada $\min(d_{ij},\,d_{ji})$: $k$ recorridos abiertos no vacíos
que cubren $n$ nodos constituyen un bosque generador de $k$ componentes. Simetrizar
hacia abajo solo puede subestimar el costo dirigido de cada arco, de modo que toda
solución factible satisface $\text{viaje} \geq \text{MSF}_k$. Es, sin embargo, una
relajación ---un bosque generador no está sujeto a la restricción de grado máximo $2$
que todo recorrido sí satisface--- y por tanto sobreestima el relleno evitable; sirve
como cota de referencia, no como objetivo alcanzable.

La distancia entre esa cota y lo que un recorrido concreto alcanza está medida. Un
experimento posterior construyó, para cada instancia, un ancla alcanzable
$\text{UB}_k$ ---los $k$ recorridos abiertos que resultan de partir un recorrido
\textit{TSP} sobre los mismos $n$ nodos por sus $k-1$ tramos más caros--- y la
contrastó contra $\text{MSF}_k$. Ese recorrido queda entre un $13$\,\% y un $29$\,\%
por encima de la cota en las instancias reales, con $13{,}3$\,\% en
\texttt{area-27-n72}, $18{,}6$\,\% en \texttt{area-26-n157}, $21{,}1$\,\% en
\texttt{area-29-n43} y a lo sumo $28{,}8$\,\% en \texttt{reference-n1607}
(\texttt{docs/experiments/tsp-achievable-anchor-20260722.md}). Dado que $\text{UB}_k$
es un recorrido construido y no óptimo, un ancla más ajustada solo podría estrechar
esa distancia, nunca ampliarla; las cifras son, por tanto, cotas superiores de la
incertidumbre del criterio. El mismo experimento advierte que el ancla no respeta
$T_{\max}$ donde $k$ es restrictivo, de modo que no constituye una solución factible
del VRP y su papel es acotar la flojedad de la cota, no sustituirla como referencia
operacional. Re-juzgadas contra el ancla, ninguna de las doce celdas de área de los
dos ciclos más recientes cambia de veredicto.

Defecto 3 --- Auto-cruces sobre cuerdas rectas.
La métrica de auto-cruces se calcula sobre cuerdas rectas entre paradas consecutivas, una
geometría que no coincide ni con la red que el \textit{solver} optimiza ni con la
polilínea real que se dibuja al censista. La correlación de Spearman entre ambas
geometrías fue de $\rho = 0{,}527$ a nivel global e invertida en el \textit{dataset} de
evaluación principal ($\rho = -0{,}575$). La prueba decisiva es que el aumento de
auto-cruces atribuido al 2-opt intra-ruta, que sobre cuerdas se multiplicaba por 2{,}9 a
5{,}3 veces, se convierte sobre la geometría real de calles en una reducción o un empate
en 104 de 108 filas comparadas. La métrica de cuerdas conserva
su valor comparativo dentro de cada ciclo, dado que todos los brazos se contrastaron con la
misma vara; las diferencias relativas permanecen válidas aunque el nivel absoluto quede
acotado por esta limitación.

Defecto 4 --- Criterio de degeneración mixto.
El criterio de rutas degeneradas original mezclaba dos fenómenos: un aislamiento
geográfico que captura el conteo de paradas (menos de cinco) y un perjuicio operativo real
que captura la duración (menos de 1\,800 segundos). Una ruta con tres paradas y
10\,485\,s de duración ocupa el 97\,\% del presupuesto diario y constituye una jornada
completa, aunque viole el conteo. En la revisión se descartó la cláusula de paradas por
indirecta, con lo que cuarenta y cinco de las ciento seis filas re-juzgables quedan
absueltas y sesenta y una persisten por duración.

\phantomsection
\subsubsection*{Alcance de las réplicas y de las barras de error}
\label{anx:alcance-replicas}

La sensibilidad al orden de entrada que la Sección~\ref{sec:metricas-calidad} mide
sobre el tiempo de desplazamiento no se limita a esa métrica. Mediciones posteriores
sobre la suite de instancias congeladas, con réplicas construidas del mismo modo,
arrojan dispersiones de otro orden de
magnitud sobre las métricas de calidad: el índice de balance de carga de una
configuración experimental, medido en 0{,}727 sobre una corrida única, resulta ser
0{,}181 $\pm$ 0{,}215 sobre cinco réplicas
(\texttt{docs/experiments/sweep-metrology-20260720.md}). La cifra no es
comparable con el 5{,}1\,\% de la Sección~\ref{sec:metricas-calidad} ---aquella es
un índice adimensional y esta dispersión del tiempo de desplazamiento---, pero fija
el orden de la sensibilidad al orden de entrada, y resulta coherente con la
diferencia de 3{,}9\,\% entre corridas distintas que tiene ese mismo origen.

Una medición posterior acota, a su vez, el alcance de esas réplicas. Sobre el
conjunto de evaluación agregado, las cinco semillas de una misma corrida arrojan
barras de error de entre $\pm 745$ y $\pm 1\,062$\,s sobre el tiempo de
desplazamiento, pero la media de una misma configuración experimental se desplazó
de $59\,971$ a $62\,751$\,s entre dos corridas distintas
(\texttt{docs/experiments/stops-penalty-sweep-20260722.md}), un salto mayor que
esa dispersión entre semillas. Las barras de error publicadas acotan, por tanto,
la varianza entre semillas dentro de una corrida, no la varianza entre corridas;
comparar medias de corridas distintas queda fuera de lo que la instrumentación
actual justifica.

\phantomsection
\subsubsection*{Efectos no medidos}
\label{anx:efectos-no-medidos}

Costo de las disyunciones de nodos.
Declarar cada árbol como una disyunción unitaria (Sección~\ref{sec:disyunciones})
añade una decisión binaria por nodo y amplía el vecindario que la búsqueda local
debe explorar en cada iteración. Como el GLS se corta por reloj de pared y no por
convergencia, cabe esperar que ese costo no se manifieste en un tiempo total mayor
sino en una calidad menor alcanzable dentro del mismo presupuesto. La magnitud del
efecto no se midió: aislarla exige una ablación con y sin disyunciones sobre el mismo
\textit{dataset}, el mismo presupuesto de \textit{solver} y las mismas semillas,
experimento que no se ejecutó. La limitación se asume porque el mecanismo no es
prescindible ---sin él un solo árbol inalcanzable deja al \textit{solver} sin solución
para todo el \textit{dataset}---, de modo que la alternativa medida no sería una
configuración operable sino una referencia de laboratorio.

Exactitud del estimador de distancias.
El perfil peatonal (\texttt{foot}) de OSRM sobre datos de \textit{OpenStreetMap}
estima el tiempo de caminata a partir de la representación de la red vial y de un
factor de velocidad constante, de modo que se aparta necesariamente de la caminata
real de terreno: no modela pendientes, semáforos, congestión peatonal, obras ni el
ritmo de una persona que se detiene a censar. La magnitud de esa desviación no está
medida en este trabajo, y no podría estarlo sin una jornada de terreno con registro de
tiempos, que queda fuera del alcance (Sección~\ref{sec:verificacion}). Cabe esperar,
eso sí, que el error sea sistemático más que aleatorio ---todas las rutas se estiman
con el mismo perfil sobre la misma red---, lo que lo hace sensible para las cifras
absolutas de duración de jornada y mucho menos sensible para las comparaciones
relativas entre configuraciones, que son las que sostienen las conclusiones del
capítulo de resultados. Contrastar la estimación contra tiempos reales de terreno es
la validación pendiente de mayor valor para el uso operativo de la plataforma.

\anexo{5}{Trazabilidad de requisitos al diseño}{Trazabilidad de requisitos al diseño}
\label{anx:trazabilidad}

La Tabla~\ref{tab:trazabilidad} remite cada requisito funcional de la
Sección~\ref{sec:requisitos-funcionales} a la sección del Capítulo~\ref{cap:diseno}
donde se documenta su realización y, cuando corresponde, al inventario de
\textit{endpoints} del \hyperref[anx:api]{Anexo 1}.

\begin{xltabular}{\textwidth}{@{}l X@{}}
  \caption{Trazabilidad de los requisitos funcionales al diseño.}
  \label{tab:trazabilidad} \\
  \toprule
  ID & Sección de diseño \\
  \midrule
  \endfirsthead
  \toprule
  ID & Sección de diseño \\
  \midrule
  \endhead
  \bottomrule
  \endlastfoot
  RF-01 & \S\ref{sec:legacy} \\
  RF-02 & \S\ref{sec:legacy} \\
  RF-03 & \S\ref{sec:legacy} \\
  RF-04 & \S\ref{sec:solver} \\
  RF-05 & \S\ref{sec:flota} \\
  RF-06 & \S\ref{sec:formulacion} \\
  RF-07 & \S\ref{sec:flota} \\
  RF-08 & \S\ref{sec:disyunciones} \\
  RF-09 & \S\ref{sec:modelo-datos} \\
  RF-10 & \S\ref{sec:ui} \\
  RF-11 & \S\ref{sec:modelo-datos} \\
  RF-12 & \S\ref{sec:modelo-datos} \\
  RF-13 & \S\ref{sec:ui} \\
  RF-14 & \S\ref{sec:ui} \\
  RF-15 & \S\ref{sec:recenso} \\
  RF-16 & \S\ref{sec:recenso} \\
  RF-17 & \S\ref{sec:recenso} \\
  RF-18 & \S\ref{sec:ui} \\
  RF-19 & \S\ref{sec:ui} \\
  RF-20 & \S\ref{sec:ui} \\
  RF-21 & \S\ref{sec:api} \\
  RF-22 & \S\ref{sec:flota} \\
  RF-23 & \S\ref{sec:ui} \\
  RF-24 & \S\ref{sec:ui} \\
  RF-25 & \S\ref{sec:legacy} \\
  RF-26 & \S\ref{sec:api} \\
\end{xltabular}

\anexo{6}{Capturas complementarias de la interfaz}{Capturas complementarias de la interfaz}
\label{anx:capturas}

% La cota de altura replica la del capítulo de diseño (capturaui).
\newcommand{\capturaanexo}[4]{%
  \begin{figure}[htbp]
    \centering
    \includegraphics[width=0.9\textwidth,height=0.75\textheight,keepaspectratio]{media/#1}
    \caption[#2]{#3}
    \label{#4}
  \end{figure}%
}

La Ilustración~\ref{fig:ui_datasets} muestra el listado de \textit{datasets} disponibles en el
sistema (\texttt{/admin/datasets}): tabla con nombre, número de árboles y fecha de importación;
los botones superiores dan acceso a las dos rutas de creación.

\capturaanexo{ui-datasets.png}{Listado de \textit{datasets} e importación de archivos}{Listado de
\textit{datasets} e importación de archivos.}{fig:ui_datasets}

Las dos rutas de creación a las que dan acceso esos botones corresponden a la
importación de archivo y a la selección sobre el mapa de fuentes legadas, ambas
descritas en la Sección~\ref{sec:ui}.

El tablero de avance del censo (\texttt{/admin/datasets/\{id\}/progress}) se muestra en la
Ilustración~\ref{fig:ui_progress}: porcentaje de árboles censados, omitidos y pendientes; mapa
coloreado por el estado de cada árbol; y desglose del avance por ruta y por censista.

\capturaanexo{ui-progress.png}{Tablero de avance del censo}{Tablero de avance del censo.}%
{fig:ui_progress}

El desglose por censista es el que permite verificar en operación el reparto de
carga que persigue el índice de balance definido en la
Sección~\ref{sec:atributos-calidad}.

\anexo{7}{Modelo de datos de las bases legadas}{Modelo de datos de las bases legadas}
\label{anx:modelo-legado}

Este anexo detalla las tablas de los dos esquemas legados que la Sección~\ref{sec:legacy} resume, y
qué columna de cada una alimenta al sistema. Ninguna de las dos bases se escribe: el acceso es
estrictamente de solo lectura (RE-03).

\begin{table}[H]
  \centering
  \caption{Esquema \texttt{arbocensus\_api\_app\_*} (\texttt{legacy\_api}, 429 árboles).}
  \small
  \begin{tabularx}{\textwidth}{@{}l r X@{}}
    \toprule
    Tabla & Filas & Columnas que usa el sistema \\
    \midrule
    \texttt{tree} & 429 &
      \texttt{id} (identidad persistente, par con \texttt{source}), \texttt{latitude},
      \texttt{longitude}, \texttt{tree\_species}. \\
    \texttt{sample} & 3\,961 &
      \texttt{real\_tree\_id} (vínculo directo al árbol), fecha, \texttt{completed} (deriva el
      estado \texttt{alive}/\texttt{unknown} de la Sección~\ref{sec:recenso}). \\
    \texttt{answer} & 12\,589 &
      \texttt{sample\_id}, \texttt{answer\_url} y \texttt{crop\_url} (10\,008 con foto; se importa
      la URL, no el archivo --- Sección~\ref{sec:recenso}). No se importan las respuestas de texto
      libre. \\
    \texttt{area} & 48 &
      \texttt{coordinates} (lista de vértices, se reconstruye como polígono), \texttt{campaign\_id}.
      27 sin arbolado, 21 no vacías que se reducen a 12 polígonos distintos por recreación entre
      campañas (Sección~\ref{sec:estado-actual}). \\
    \texttt{campaign} & 10 &
      Nombre, usado solo como metadato de exploración; no se importa como entidad propia. \\
    \texttt{step} & --- &
      Definición de preguntas por campaña. No se importa: el sistema no reconstruye el cuestionario
      histórico, solo el estado derivado de \texttt{completed}. \\
    \bottomrule
  \end{tabularx}
\end{table}

\begin{table}[H]
  \centering
  \caption{Esquema \texttt{arbocensus\_app\_*} (\texttt{legacy\_app}, 12\,941 árboles con posición
    utilizable).}
  \small
  \begin{tabularx}{\textwidth}{@{}l r X@{}}
    \toprule
    Tabla & Filas & Columnas que usa el sistema \\
    \midrule
    \texttt{sample} & 17\,208 &
      Código QR (identidad persistente, par con \texttt{source}), latitud y longitud de cada lectura
      de GPS. La posición del árbol se reconstruye como la mediana de las lecturas agrupadas por QR
      (Sección~\ref{sec:legacy}). \\
    \texttt{answer} & --- &
      \texttt{answer\_url} (83\,614 URL \texttt{https}; se importan solo las que comienzan con
      \texttt{http}, descartando rutas locales del dispositivo). \texttt{completed} deriva el mismo
      estado que en \texttt{legacy\_api}. \\
    \bottomrule
  \end{tabularx}
\end{table}

La diferencia estructural entre ambos esquemas es la que obliga al doble camino de vinculación
descrito en la Sección~\ref{sec:recenso}: \texttt{legacy\_api} referencia el árbol directamente desde
la observación (\texttt{sample.real\_tree\_id}), mientras que \texttt{legacy\_app} no tiene una tabla
de árboles propia y el árbol se reconstruye a partir del código QR que agrupa sus muestras. La
identidad persistente que \texttt{Tree} conserva en toda importación es, en ambos casos, el par
(\texttt{source}, \texttt{external\_id}) y no el identificador de fila de la base legada.
